\documentclass[11pt,a4paper]{article}

\usepackage[utf8]{inputenc}
\usepackage{amsmath,amssymb,amsthm}
\usepackage[margin=1in]{geometry}
\usepackage{hyperref}
\usepackage{enumitem}
\usepackage{booktabs}

\newtheorem{theorem}{Theorem}[section]
\newtheorem{lemma}[theorem]{Lemma}
\newtheorem{proposition}[theorem]{Proposition}
\newtheorem{corollary}[theorem]{Corollary}
\newtheorem{remark}[theorem]{Remark}
\theoremstyle{definition}
\newtheorem{definition}[theorem]{Definition}
\newtheorem{hypothesis}[theorem]{Hypothesis}

\newcommand{\norm}[1]{\lVert #1 \rVert}
\newcommand{\abs}[1]{\lvert #1 \rvert}
\newcommand{\R}{\mathbb{R}}
\newcommand{\N}{\mathbb{N}}
\newcommand{\Icc}[2]{[#1,\, #2]}
\newcommand{\Ico}[2]{[#1,\, #2)}
\newcommand{\dist}{\mathrm{dist}}

\title{Blueprint: Convergence of the Euler Method for ODEs}
\author{Formalization Blueprint}
\date{}

\begin{document}
\maketitle

\tableofcontents
\newpage

%% ============================================================
\section{Setup and Hypotheses}\label{sec:setup}
%% ============================================================

Let $E$ be a real Banach space (i.e., a complete normed vector space over~$\R$).
Fix real numbers $a < b$ and an initial value $x_0 \in E$.
Let $v \colon \R \times E \to E$ be a vector field.

We impose the following three hypotheses on~$v$.

\begin{hypothesis}[Lipschitz in space]\label{hyp:lip-x}
There exists a constant $K \ge 0$ such that for every $t \in \Icc{a}{b}$
and every $x, y \in E$,
\[
  \norm{v(t,x) - v(t,y)} \;\le\; K\,\norm{x - y}.
\]
\end{hypothesis}

\begin{hypothesis}[Lipschitz in time]\label{hyp:lip-t}
There exists a constant $L \ge 0$ such that for every $t, s \in \Icc{a}{b}$
and every $x \in E$,
\[
  \norm{v(t,x) - v(s,x)} \;\le\; L\,\abs{t - s}.
\]
\end{hypothesis}

\begin{hypothesis}[Boundedness]\label{hyp:bdd}
There exists a constant $M \ge 0$ such that for every $t \in \Icc{a}{b}$
and every $x \in E$,
\[
  \norm{v(t,x)} \;\le\; M.
\]
\end{hypothesis}

\begin{remark}
Hypothesis~\ref{hyp:bdd} is a strong global assumption.
In practice one works on a ball $\{x : \norm{x - x_0} \le R\}$
with $R > M(b-a)$, and verifies that all approximations remain
in this ball. The global formulation simplifies the proof and
suffices for a first formalization.
\end{remark}

\begin{remark}\label{rem:continuity}
Hypotheses~\ref{hyp:lip-x} and~\ref{hyp:lip-t} together imply
that the map $(t, x) \mapsto v(t, x)$ is continuous on
$\Icc{a}{b} \times E$. Indeed, for any $(t,x)$ and $(s,y)$
in $\Icc{a}{b} \times E$,
\[
  \norm{v(t,x) - v(s,y)}
  \;\le\; \norm{v(t,x) - v(s,x)} + \norm{v(s,x) - v(s,y)}
  \;\le\; L\,\abs{t - s} + K\,\norm{x - y}.
\]
\end{remark}


%% ============================================================
\section{The Euler Method}\label{sec:euler-def}
%% ============================================================

Throughout this section, fix a positive integer $N \ge 1$.
Define the \emph{step size} and \emph{grid points} by
\[
  h \;=\; \frac{b - a}{N}, \qquad
  t_k \;=\; a + k\,h, \qquad k = 0, 1, \ldots, N.
\]
Note that $t_0 = a$, $t_N = b$, and $t_{k+1} - t_k = h$ for all~$k$.

\begin{definition}[Euler grid values]\label{def:euler-grid}
Define the sequence $x_0, x_1, \ldots, x_N \in E$ recursively by:
\begin{align}
  x_0 &\;=\; x_0, \label{eq:euler-init}\\
  x_{k+1} &\;=\; x_k + h \cdot v(t_k,\, x_k), \qquad k = 0, 1, \ldots, N-1.
  \label{eq:euler-step}
\end{align}
\end{definition}

\begin{definition}[Euler interpolant]\label{def:euler-interp}
Define the \emph{piecewise-linear Euler interpolant}
$f_h \colon \Icc{a}{b} \to E$ as follows. For $t \in \Icc{t_k}{t_{k+1}}$
(with $0 \le k \le N-1$), set
\begin{equation}\label{eq:interp-def}
  f_h(t) \;=\; x_k + (t - t_k)\,v(t_k,\, x_k).
\end{equation}
Equivalently, $f_h$ is the unique continuous piecewise-affine
function on $\Icc{a}{b}$ that is affine on each subinterval
$\Icc{t_k}{t_{k+1}}$ and satisfies $f_h(t_k) = x_k$ for each~$k$.
\end{definition}

\begin{remark}\label{rem:index-function}
Given $t \in \Ico{a}{b}$, the index $k = k(t)$ is uniquely
determined by $t \in \Ico{t_k}{t_{k+1}}$.
Explicitly,
\[
  k(t) \;=\; \Big\lfloor \frac{t - a}{h} \Big\rfloor
\]
where $\lfloor \cdot \rfloor$ denotes the floor function.
For $t = b$ we set $k(b) = N - 1$, so that $b \in \Icc{t_{N-1}}{t_N}$.
\end{remark}


%% ============================================================
\section{Properties of the Euler Interpolant}\label{sec:interp-props}
%% ============================================================

\begin{lemma}[Grid-point values]\label{lem:grid-values}
For each $k \in \{0, 1, \ldots, N\}$,
\[
  f_h(t_k) \;=\; x_k.
\]
\end{lemma}

\begin{proof}
For $k = 0$: by~\eqref{eq:interp-def} with $t = t_0$,
$f_h(t_0) = x_0 + 0 \cdot v(t_0, x_0) = x_0$.

For $k \ge 1$: at $t = t_k$, using the piece on
$\Icc{t_{k-1}}{t_k}$, we get
\begin{align*}
  f_h(t_k)
  &= x_{k-1} + (t_k - t_{k-1})\,v(t_{k-1},\, x_{k-1}) \\
  &= x_{k-1} + h\,v(t_{k-1},\, x_{k-1}) \\
  &= x_k,
\end{align*}
where the last equality is the Euler recurrence~\eqref{eq:euler-step}.
This also equals $x_k + (t_k - t_k)\,v(t_k, x_k) = x_k$
from the piece on $\Icc{t_k}{t_{k+1}}$ (when $k < N$),
confirming that the two pieces agree at~$t_k$.
\end{proof}


\begin{lemma}[Continuity]\label{lem:continuity}
The function $f_h$ is continuous on $\Icc{a}{b}$.
\end{lemma}

\begin{proof}
On each closed subinterval $\Icc{t_k}{t_{k+1}}$, the map
$t \mapsto x_k + (t - t_k)\,v(t_k, x_k)$ is affine (hence
continuous). By Lemma~\ref{lem:grid-values}, the value at
$t_k$ from the piece on $\Icc{t_{k-1}}{t_k}$ equals $x_k$,
and the value at $t_k$ from the piece on $\Icc{t_k}{t_{k+1}}$
also equals~$x_k$. So the pieces agree at every grid point, and
$f_h$ is continuous on $\Icc{a}{b}$.
\end{proof}


\begin{lemma}[Right derivative]\label{lem:right-deriv}
For every $t \in \Ico{a}{b}$, let $k = k(t)$ be the index such
that $t \in \Ico{t_k}{t_{k+1}}$. Then $f_h$ has a right derivative
at~$t$, equal to $v(t_k, x_k)$. More precisely, the function $f_h$
is differentiable from the right at $t$ with
\begin{equation}\label{eq:right-deriv}
  \lim_{s \to t^+} \frac{f_h(s) - f_h(t)}{s - t}
  \;=\; v(t_k,\, x_k).
\end{equation}
\end{lemma}

\begin{proof}
\textbf{Case 1: $t \in (t_k, t_{k+1})$.}
Then for $s > t$ sufficiently close to~$t$, we have
$s \in (t_k, t_{k+1})$ as well, so
\[
  f_h(s) - f_h(t) = \bigl[x_k + (s - t_k)\,v(t_k, x_k)\bigr]
                   - \bigl[x_k + (t - t_k)\,v(t_k, x_k)\bigr]
                   = (s - t)\,v(t_k, x_k).
\]
Dividing by $s - t > 0$ gives $v(t_k, x_k)$.

\textbf{Case 2: $t = t_k$ for some $k < N$.}
For $s \in (t_k, t_{k+1})$,
\[
  f_h(s) - f_h(t_k)
  = x_k + (s - t_k)\,v(t_k, x_k) - x_k
  = (s - t_k)\,v(t_k, x_k).
\]
Dividing by $s - t_k > 0$ gives $v(t_k, x_k)$.
\end{proof}


\begin{lemma}[Deviation from grid point]\label{lem:deviation}
For $t \in \Icc{t_k}{t_{k+1}}$,
\begin{equation}\label{eq:deviation}
  \norm{f_h(t) - x_k} \;=\; \abs{t - t_k}\,\norm{v(t_k, x_k)}
  \;\le\; h\,M.
\end{equation}
\end{lemma}

\begin{proof}
By~\eqref{eq:interp-def},
$f_h(t) - x_k = (t - t_k)\,v(t_k, x_k)$, so
\[
  \norm{f_h(t) - x_k}
  = \abs{t - t_k}\,\norm{v(t_k, x_k)}
  \le h\,M,
\]
where we used $\abs{t - t_k} \le h$ and Hypothesis~\ref{hyp:bdd}.
\end{proof}


\begin{lemma}[Defect bound]\label{lem:defect}
Define the \emph{defect} of the Euler interpolant at $t \in \Ico{a}{b}$ as
\[
  d(t) \;=\; \norm{f_h'(t) - v\bigl(t,\, f_h(t)\bigr)},
\]
where $f_h'(t) = v(t_k, x_k)$ is the right derivative from
Lemma~\ref{lem:right-deriv}, and $k = k(t)$. Then
\begin{equation}\label{eq:defect-bound}
  d(t) \;\le\; (L + K\,M)\,h
\end{equation}
for all $t \in \Ico{a}{b}$.
\end{lemma}

\begin{proof}
Let $t \in \Ico{t_k}{t_{k+1}}$. By the triangle inequality,
\begin{equation}\label{eq:defect-split}
  d(t) \;=\; \norm{v(t_k, x_k) - v(t, f_h(t))}
  \;\le\; \underbrace{\norm{v(t_k, x_k) - v(t, x_k)}}_{I_1}
       +  \underbrace{\norm{v(t, x_k) - v(t, f_h(t))}}_{I_2}.
\end{equation}

\textbf{Bound on $I_1$.}
By Hypothesis~\ref{hyp:lip-t},
\begin{equation}\label{eq:I1-bound}
  I_1 = \norm{v(t_k, x_k) - v(t, x_k)}
  \le L\,\abs{t_k - t}
  \le L\,h.
\end{equation}

\textbf{Bound on $I_2$.}
By Hypothesis~\ref{hyp:lip-x},
\[
  I_2 = \norm{v(t, x_k) - v(t, f_h(t))}
  \le K\,\norm{x_k - f_h(t)}.
\]
By Lemma~\ref{lem:deviation}, $\norm{x_k - f_h(t)} \le h\,M$, so
\begin{equation}\label{eq:I2-bound}
  I_2 \;\le\; K\,M\,h.
\end{equation}

Combining~\eqref{eq:I1-bound} and~\eqref{eq:I2-bound}:
\[
  d(t) \;\le\; L\,h + K\,M\,h \;=\; (L + K\,M)\,h.
  \qedhere
\]
\end{proof}


%% ============================================================
\section{Existence of the True Solution}\label{sec:existence}
%% ============================================================

\begin{lemma}[Existence and regularity of the ODE solution]\label{lem:true-solution}
Under Hypotheses~\ref{hyp:lip-x}--\ref{hyp:bdd}, there exists
a function $\alpha \colon \R \to E$ satisfying:
\begin{enumerate}[label=(\roman*)]
  \item\label{it:init} $\alpha(a) = x_0$.
  \item\label{it:cont} $\alpha$ is continuous on $\Icc{a}{b}$.
  \item\label{it:deriv} For every $t \in \Ico{a}{b}$,
    $\alpha$ has a right derivative at $t$ equal to $v(t, \alpha(t))$:
    \[
      \lim_{s \to t^+} \frac{\alpha(s) - \alpha(t)}{s - t}
      \;=\; v\bigl(t,\, \alpha(t)\bigr).
    \]
    Equivalently, $\alpha$ satisfies the ODE
    $\alpha'(t) = v(t, \alpha(t))$ on $\Ico{a}{b}$.
  \item\label{it:integral} For every $t \in \Icc{a}{b}$,
    \begin{equation}\label{eq:integral-form}
      \alpha(t) \;=\; x_0 + \int_a^t v\bigl(s,\, \alpha(s)\bigr)\,ds.
    \end{equation}
  \item\label{it:exact-defect}
    The defect of $\alpha$ as an approximate solution is zero:
    \[
      \norm{\alpha'(t) - v\bigl(t,\, \alpha(t)\bigr)} = 0
      \quad \text{for all } t \in \Ico{a}{b}.
    \]
\end{enumerate}
\end{lemma}

\begin{proof}
We verify the hypotheses of the Picard--Lindel\"of theorem.
Set the ball radius $a_{\mathrm{ball}} = M(b - a)$ and $r = 0$.
Then:
\begin{itemize}
  \item By Hypothesis~\ref{hyp:lip-x}, for each
    $t \in \Icc{a}{b}$, the map $v(t, \cdot)$ is $K$-Lipschitz on
    the closed ball $\overline{B}(x_0, a_{\mathrm{ball}})$.
  \item By Remark~\ref{rem:continuity}, for each
    $x \in \overline{B}(x_0, a_{\mathrm{ball}})$, the map
    $t \mapsto v(t, x)$ is continuous on~$\Icc{a}{b}$.
  \item By Hypothesis~\ref{hyp:bdd},
    $\norm{v(t, x)} \le M$ for $t \in \Icc{a}{b}$,
    $x \in \overline{B}(x_0, a_{\mathrm{ball}})$.
  \item The time--interval condition:
    $M \cdot (b - a) \le a_{\mathrm{ball}} - r = M(b - a) - 0 = M(b - a)$.
    This holds with equality.
\end{itemize}
The Picard--Lindel\"of theorem yields
$\alpha \colon \R \to E$ satisfying~\ref{it:init}--\ref{it:deriv}.

Property~\ref{it:integral} follows from~\ref{it:deriv} by the
fundamental theorem of calculus: since $\alpha$ is continuous on
$\Icc{a}{b}$ and has right derivative $v(t, \alpha(t))$ on
$\Ico{a}{b}$, and since $t \mapsto v(t, \alpha(t))$ is continuous
(by Remark~\ref{rem:continuity} and the continuity of~$\alpha$),
we have $\alpha(t) - \alpha(a) = \int_a^t v(s, \alpha(s))\,ds$.

Property~\ref{it:exact-defect} is immediate from~\ref{it:deriv}.
\end{proof}


\begin{lemma}[Uniqueness]\label{lem:uniqueness}
Under Hypothesis~\ref{hyp:lip-x}, the solution $\alpha$ from
Lemma~\ref{lem:true-solution} is unique on $\Icc{a}{b}$.
That is, if $\beta \colon \R \to E$ also satisfies $\beta(a) = x_0$
and $\beta'(t) = v(t, \beta(t))$ for all $t \in \Ico{a}{b}$,
with $\beta$ continuous on $\Icc{a}{b}$, then
$\beta(t) = \alpha(t)$ for all $t \in \Icc{a}{b}$.
\end{lemma}

\begin{proof}
This follows from the Gr\"onwall-based uniqueness theorem
for ODEs with Lipschitz right-hand side.
\end{proof}


%% ============================================================
\section{The Gr\"onwall Inequality for Approximate Solutions}
\label{sec:gronwall}
%% ============================================================

The following theorem is the key tool that converts the local
defect bound (Lemma~\ref{lem:defect}) into a global error bound.

\begin{theorem}[Gr\"onwall bound for approximate solutions]
\label{thm:gronwall-approx}
Let $f, g \colon \Icc{a}{b} \to E$ be continuous functions that
possess right derivatives $f'(t)$ and $g'(t)$ for all
$t \in \Ico{a}{b}$. Suppose that $v(t, \cdot)$ is $K$-Lipschitz
for every $t$. Suppose there exist constants
$\varepsilon_f, \varepsilon_g, \delta \ge 0$ such that:
\begin{enumerate}[label=(\alph*)]
  \item $\norm{f'(t) - v(t, f(t))} \le \varepsilon_f$
    for all $t \in \Ico{a}{b}$,
  \item $\norm{g'(t) - v(t, g(t))} \le \varepsilon_g$
    for all $t \in \Ico{a}{b}$,
  \item $\norm{f(a) - g(a)} \le \delta$.
\end{enumerate}
Then for all $t \in \Icc{a}{b}$,
\begin{equation}\label{eq:gronwall-bound}
  \norm{f(t) - g(t)}
  \;\le\; G\bigl(\delta,\; K,\; \varepsilon_f + \varepsilon_g,\;
                    t - a\bigr),
\end{equation}
where the \emph{Gr\"onwall bound} $G$ is defined by
\begin{equation}\label{eq:gronwall-def}
  G(\delta, K, \varepsilon, x) \;=\;
  \begin{cases}
    \delta\,e^{Kx} + \dfrac{\varepsilon}{K}\,\bigl(e^{Kx} - 1\bigr)
    & \text{if } K \neq 0, \\[6pt]
    \delta + \varepsilon\,x
    & \text{if } K = 0.
  \end{cases}
\end{equation}
\end{theorem}

\begin{proof}
This is a standard consequence of the Gr\"onwall inequality.
The proof proceeds by bounding the right derivative of
$t \mapsto \norm{f(t) - g(t)}$ and applying the differential
Gr\"onwall lemma.

Define $e(t) = f(t) - g(t)$. Then $e$ is continuous on $\Icc{a}{b}$,
and for $t \in \Ico{a}{b}$, the right derivative of $e$ is
$e'(t) = f'(t) - g'(t)$. We estimate:
\begin{align*}
  \norm{e'(t)}
  &= \norm{f'(t) - g'(t)} \\
  &\le \norm{f'(t) - v(t, f(t))} + \norm{v(t, f(t)) - v(t, g(t))}
     + \norm{v(t, g(t)) - g'(t)} \\
  &\le \varepsilon_f + K\,\norm{f(t) - g(t)} + \varepsilon_g \\
  &= K\,\norm{e(t)} + (\varepsilon_f + \varepsilon_g).
\end{align*}
Since $\norm{e(\cdot)}$ is continuous and its right derivative
$\norm{e'(\cdot)}$ satisfies
$\norm{e'(t)} \le K\,\norm{e(t)} + (\varepsilon_f + \varepsilon_g)$,
the scalar Gr\"onwall inequality (applied to $\norm{e(\cdot)}$)
gives~\eqref{eq:gronwall-bound}.
\end{proof}


%% ============================================================
\section{Main Convergence Theorem}\label{sec:convergence}
%% ============================================================

\begin{theorem}[Convergence of the Euler method]
\label{thm:main}
Let $E$ be a real Banach space, let $a < b$ be real numbers,
let $x_0 \in E$, and let $v \colon \R \times E \to E$ satisfy
Hypotheses~\ref{hyp:lip-x}--\ref{hyp:bdd} with constants $K, L, M$.

Let $\alpha$ be the unique solution to the initial-value problem
$\alpha'(t) = v(t, \alpha(t))$, $\alpha(a) = x_0$, on $\Icc{a}{b}$
(which exists by Lemma~\ref{lem:true-solution}).

For each positive integer $N$, let $f_{(b-a)/N}$ be the Euler
interpolant from Definition~\ref{def:euler-interp} with step size
$h = (b-a)/N$.

Then for all $t \in \Icc{a}{b}$,
\begin{equation}\label{eq:main-bound}
  \norm{f_h(t) - \alpha(t)}
  \;\le\; G\bigl(0,\; K,\; (L + KM)\,h,\; t - a\bigr).
\end{equation}
Explicitly:
\begin{equation}\label{eq:main-explicit}
  \norm{f_h(t) - \alpha(t)}
  \;\le\;
  \begin{cases}
    \dfrac{(L + KM)\,h}{K}\,
    \bigl(e^{K(t-a)} - 1\bigr)
    & \text{if } K > 0, \\[6pt]
    L\,h\,(t - a)
    & \text{if } K = 0.
  \end{cases}
\end{equation}
In particular, uniformly on $\Icc{a}{b}$,
\begin{equation}\label{eq:main-uniform}
  \sup_{t \in \Icc{a}{b}} \norm{f_h(t) - \alpha(t)}
  \;\le\; C\,h,
\end{equation}
where the constant $C$ depends only on $K$, $L$, $M$, and $b - a$:
\begin{equation}\label{eq:C-def}
  C \;=\;
  \begin{cases}
    \dfrac{L + KM}{K}\,\bigl(e^{K(b-a)} - 1\bigr)
    & \text{if } K > 0, \\[6pt]
    L\,(b-a)
    & \text{if } K = 0.
  \end{cases}
\end{equation}
Consequently, as $N \to \infty$ (equivalently, $h \to 0$),
the Euler interpolants $f_h$ converge uniformly to $\alpha$ on
$\Icc{a}{b}$.
\end{theorem}

\begin{proof}
We apply Theorem~\ref{thm:gronwall-approx} with:
\begin{itemize}
  \item $f = f_h$ (the Euler interpolant),
  \item $g = \alpha$ (the true solution),
  \item $\varepsilon_f = (L + KM)\,h$ (from Lemma~\ref{lem:defect}),
  \item $\varepsilon_g = 0$ (from Lemma~\ref{lem:true-solution}\ref{it:exact-defect}),
  \item $\delta = 0$ (since $f_h(a) = x_0 = \alpha(a)$).
\end{itemize}

We verify the hypotheses of Theorem~\ref{thm:gronwall-approx}:
\begin{enumerate}[label=(\alph*)]
  \item \textbf{$f_h$ is continuous on $\Icc{a}{b}$.}
    This is Lemma~\ref{lem:continuity}.

  \item \textbf{$f_h$ has a right derivative at every $t \in \Ico{a}{b}$.}
    This is Lemma~\ref{lem:right-deriv}. The right derivative
    is $f_h'(t) = v(t_k, x_k)$ where $k = k(t)$.

  \item \textbf{$\alpha$ is continuous on $\Icc{a}{b}$.}
    This is Lemma~\ref{lem:true-solution}\ref{it:cont}.

  \item \textbf{$\alpha$ has a right derivative at every $t \in \Ico{a}{b}$.}
    This is Lemma~\ref{lem:true-solution}\ref{it:deriv}, with
    right derivative $\alpha'(t) = v(t, \alpha(t))$.

  \item \textbf{Defect of $f_h$:}
    By Lemma~\ref{lem:defect},
    $\norm{f_h'(t) - v(t, f_h(t))} \le (L + KM)\,h$ for all
    $t \in \Ico{a}{b}$.

  \item \textbf{Defect of $\alpha$:}
    $\norm{\alpha'(t) - v(t, \alpha(t))} = 0 \le 0$ for all
    $t \in \Ico{a}{b}$.

  \item \textbf{Initial condition:}
    $\norm{f_h(a) - \alpha(a)} = \norm{x_0 - x_0} = 0$.
\end{enumerate}

Theorem~\ref{thm:gronwall-approx} now gives, for all $t \in \Icc{a}{b}$:
\[
  \norm{f_h(t) - \alpha(t)}
  \;\le\; G\bigl(0,\; K,\; (L + KM)\,h + 0,\; t - a\bigr).
\]

\textbf{Case $K > 0$.}
By~\eqref{eq:gronwall-def} with $\delta = 0$,
\[
  G(0, K, \varepsilon, x)
  = 0 \cdot e^{Kx} + \frac{\varepsilon}{K}\,(e^{Kx} - 1)
  = \frac{\varepsilon}{K}\,(e^{Kx} - 1).
\]
With $\varepsilon = (L + KM)\,h$ and $x = t - a \le b - a$:
\[
  \norm{f_h(t) - \alpha(t)}
  \le \frac{(L + KM)\,h}{K}\,(e^{K(t-a)} - 1)
  \le \frac{(L + KM)}{K}\,(e^{K(b-a)} - 1)\,h.
\]

\textbf{Case $K = 0$.}
By~\eqref{eq:gronwall-def},
$G(0, 0, \varepsilon, x) = 0 + \varepsilon\,x = \varepsilon\,x$.
Since $K = 0$, we have $\varepsilon = (L + 0)\,h = L\,h$, so
\[
  \norm{f_h(t) - \alpha(t)}
  \le L\,h\,(t - a)
  \le L\,(b - a)\,h.
\]

In both cases, $\norm{f_h(t) - \alpha(t)} \le C\,h$ with $C$
as in~\eqref{eq:C-def}. Since $h = (b-a)/N \to 0$ as $N \to \infty$,
the convergence is established.
\end{proof}


%% ============================================================
\section{Dependency Graph}\label{sec:dependency}
%% ============================================================

The logical dependencies among the results are:

\medskip
\begin{center}
\begin{tabular}{@{}lcl@{}}
\toprule
\textbf{Result} & \textbf{Depends on} & \textbf{Section} \\
\midrule
Lemma~\ref{lem:grid-values} (grid values)
  & Def.~\ref{def:euler-grid}, Def.~\ref{def:euler-interp}
  & \S\ref{sec:interp-props} \\
Lemma~\ref{lem:continuity} (continuity)
  & Lemma~\ref{lem:grid-values}
  & \S\ref{sec:interp-props} \\
Lemma~\ref{lem:right-deriv} (right derivative)
  & Def.~\ref{def:euler-interp}
  & \S\ref{sec:interp-props} \\
Lemma~\ref{lem:deviation} (deviation)
  & Def.~\ref{def:euler-interp}, Hyp.~\ref{hyp:bdd}
  & \S\ref{sec:interp-props} \\
Lemma~\ref{lem:defect} (defect bound)
  & Lemma~\ref{lem:deviation}, Hyp.~\ref{hyp:lip-x},
    Hyp.~\ref{hyp:lip-t}
  & \S\ref{sec:interp-props} \\
Lemma~\ref{lem:true-solution} (existence)
  & Hyp.~\ref{hyp:lip-x}--\ref{hyp:bdd},
    Picard--Lindel\"of (external)
  & \S\ref{sec:existence} \\
Theorem~\ref{thm:gronwall-approx} (Gr\"onwall)
  & Hyp.~\ref{hyp:lip-x}, Gr\"onwall inequality (external)
  & \S\ref{sec:gronwall} \\
\textbf{Theorem~\ref{thm:main}} (main)
  & Lemmas~\ref{lem:continuity},~\ref{lem:right-deriv},
    \ref{lem:defect},~\ref{lem:true-solution},
    Thm.~\ref{thm:gronwall-approx}
  & \S\ref{sec:convergence} \\
\bottomrule
\end{tabular}
\end{center}


%% ============================================================
\section{Existing Mathlib Results}\label{sec:mathlib}
%% ============================================================

This section catalogs the Mathlib theorems and definitions that
should be used (not re-proved) in the formalization.

\subsection{ODE Existence: Picard--Lindel\"of}
\label{ssec:picard}

File: \texttt{Mathlib.Analysis.ODE.PicardLindelof}

\begin{itemize}
  \item \texttt{IsPicardLindelof} --- Prop structure bundling the
    hypotheses: Lipschitz in~$x$ on a closed ball, continuous in~$t$,
    norm bound, time--interval condition.

  \item \texttt{IsPicardLindelof.exists\_eq\_forall\_mem\_Icc\_hasDerivWithinAt}
    --- Existence of a solution $\alpha$ with $\alpha(t_0) = x$ and
    $\mathrm{HasDerivWithinAt}\;\alpha\;(f\;t\;(\alpha\;t))\;
    (\mathrm{Icc}\;t_{\min}\;t_{\max})\;t$
    for all $t \in \Icc{t_{\min}}{t_{\max}}$.
    \emph{Use this for Lemma~\ref{lem:true-solution}.}

  \item \texttt{IsPicardLindelof.of\_time\_independent} ---
    Simplified constructor for autonomous vector fields $v(x)$.

  \item \texttt{ODE.picard\_apply} ---
    The Picard iteration: $\mathrm{picard}\;f\;t_0\;x_0\;\alpha\;t
    = x_0 + \int_{t_0}^{t} f(\tau, \alpha(\tau))\,d\tau$.
\end{itemize}


\subsection{Gr\"onwall Inequality}
\label{ssec:gronwall}

File: \texttt{Mathlib.Analysis.ODE.Gronwall}

\begin{itemize}
  \item \texttt{gronwallBound} --- The function
    $G(\delta, K, \varepsilon, x)$ from~\eqref{eq:gronwall-def}.

  \item \texttt{gronwallBound\_x0} ---
    $G(\delta, K, \varepsilon, 0) = \delta$.

  \item \texttt{gronwallBound\_of\_K\_ne\_0} ---
    The explicit formula when $K \neq 0$.

  \item \texttt{gronwallBound\_K0} ---
    The explicit formula when $K = 0$: $G(\delta, 0, \varepsilon, x) = \delta + \varepsilon\,x$.

  \item \texttt{gronwallBound\_epsilon0} ---
    $G(\delta, K, 0, x) = \delta\,e^{Kx}$.

  \item \texttt{dist\_le\_of\_approx\_trajectories\_ODE} ---
    \textbf{This is the central external theorem.}
    Given two approximate solutions $f$ and $g$ with defects
    $\varepsilon_f$ and $\varepsilon_g$ and initial distance $\delta$,
    concludes $\dist(f(t), g(t)) \le G(\delta, K,
    \varepsilon_f + \varepsilon_g, t - a)$.
    Hypotheses:
    \begin{itemize}
      \item \texttt{hv : $\forall$ t, LipschitzWith K (v t)}
      \item \texttt{hf : ContinuousOn f (Icc a b)}
      \item \texttt{hf' : $\forall$ t $\in$ Ico a b,
        HasDerivWithinAt f (f' t) (Ici t) t}
      \item \texttt{f\_bound : $\forall$ t $\in$ Ico a b,
        dist (f' t) (v t (f t)) $\le$ $\varepsilon_f$}
      \item (same for $g$, $g'$, $\varepsilon_g$)
      \item \texttt{ha : dist (f a) (g a) $\le$ $\delta$}
    \end{itemize}
    \emph{Use this for the main step of Theorem~\ref{thm:main}.}

  \item \texttt{ODE\_solution\_unique} ---
    Uniqueness of ODE solutions.
    \emph{Use this for Lemma~\ref{lem:uniqueness}.}
\end{itemize}


\subsection{Interval Integrals and Fundamental Theorem of Calculus}
\label{ssec:ftc}

File: \texttt{Mathlib.MeasureTheory.Integral.IntervalIntegral}
and related files.

\begin{itemize}
  \item \texttt{intervalIntegral.integral\_eq\_sub\_of\_hasDerivAt}
    --- FTC-2: if $f$ has derivative $f'$ on $\Icc{a}{b}$ and $f'$
    is integrable, then $\int_a^b f'(t)\,dt = f(b) - f(a)$.
    \emph{Used to establish the integral form~\eqref{eq:integral-form}
    of the ODE solution.}

  \item \texttt{ContinuousOn.intervalIntegrable} --- A continuous
    function on $\Icc{a}{b}$ is interval-integrable.
\end{itemize}


\subsection{Lipschitz Functions}
\label{ssec:lipschitz}

Files: \texttt{Mathlib.Topology.MetricSpace.Lipschitz},
\texttt{Mathlib.Topology.MetricSpace.Holder}

\begin{itemize}
  \item \texttt{LipschitzWith} --- The predicate that a function
    is $K$-Lipschitz: $\dist(f(x), f(y)) \le K \cdot \dist(x, y)$.

  \item \texttt{LipschitzWith.dist\_le\_mul} ---
    If \texttt{LipschitzWith K f}, then
    $\dist(f(x), f(y)) \le K \cdot \dist(x, y)$.
    \emph{Use this to bound $\norm{v(t,x) - v(t,y)} \le K\,\norm{x-y}$
    in Lemma~\ref{lem:defect}.}

  \item \texttt{LipschitzWith.continuous} --- A Lipschitz function
    is continuous.
\end{itemize}


\subsection{Norms, Distances, and Basic Estimates}
\label{ssec:norms}

\begin{itemize}
  \item \texttt{norm\_smul} ---
    $\norm{c \cdot x} = \abs{c} \cdot \norm{x}$ for $c \in \R$,
    $x \in E$.
    \emph{Used in Lemma~\ref{lem:deviation}:
    $\norm{(t - t_k) \cdot v(t_k, x_k)} = \abs{t - t_k} \cdot \norm{v(t_k, x_k)}$.}

  \item \texttt{dist\_comm}, \texttt{dist\_eq\_norm} ---
    $\dist(x, y) = \norm{x - y}$.

  \item \texttt{norm\_sub\_le} ---
    Triangle inequality: $\norm{x - z} \le \norm{x - y} + \norm{y - z}$.
    \emph{Used in Lemma~\ref{lem:defect}.}
\end{itemize}


\subsection{HasDerivWithinAt and Differentiability}
\label{ssec:deriv}

File: \texttt{Mathlib.Analysis.Calculus.Deriv.Basic}

\begin{itemize}
  \item \texttt{HasDerivWithinAt} ---
    The predicate that the derivative of $f$ within a set $s$ at
    a point $x$ is $y$. For the Euler interpolant and the Gronwall
    theorem, the relevant set is $\mathrm{Ici}\;t = [t, \infty)$
    (right derivative).

  \item \texttt{hasDerivWithinAt\_const}, \texttt{HasDerivWithinAt.add},
    \texttt{HasDerivWithinAt.smul} ---
    Rules for computing derivatives of constants, sums, and scalar
    multiples. \emph{Used to show that the affine map
    $t \mapsto x_k + (t - t_k) \cdot c$ has derivative $c$
    within $[t, \infty)$.}

  \item \texttt{hasDerivAt\_id'} --- The identity function has
    derivative~$1$.
\end{itemize}


\subsection{Continuous Functions on Intervals}
\label{ssec:continuous}

\begin{itemize}
  \item \texttt{ContinuousOn} --- The predicate that a function
    is continuous on a set.

  \item \texttt{ContinuousOn.mono} --- If $f$ is continuous on $s$
    and $s' \subseteq s$, then $f$ is continuous on $s'$.

  \item Piecewise continuity: to show $f_h$ is continuous on
    $\Icc{a}{b}$, one approach is to show it is continuous on
    each $\Icc{t_k}{t_{k+1}}$ and that the values agree at
    grid points. Alternatively, use
    \texttt{ContinuousOn.if} or construct $f_h$ as a
    continuous map directly.
\end{itemize}


\subsection{Floor Function and Interval Arithmetic}
\label{ssec:floor}

\begin{itemize}
  \item \texttt{Nat.floor} / \texttt{Int.floor} ---
    The floor function. Used to define $k(t)$ in
    Remark~\ref{rem:index-function}.

  \item \texttt{Nat.floor\_le}, \texttt{Nat.lt\_floor\_add\_one} ---
    Properties: $\lfloor x \rfloor \le x < \lfloor x \rfloor + 1$.
\end{itemize}


%% ============================================================
\section{Implementation Notes}\label{sec:implementation}
%% ============================================================

\begin{enumerate}
  \item \textbf{Defining the interpolant.}
    The cleanest Lean definition of $f_h$ may be as a single function
    using the floor function to determine the subinterval index,
    rather than as a literal piecewise definition using
    \texttt{Set.piecewise}. Specifically:
    \begin{verbatim}
    noncomputable def eulerInterp (v : R -> E -> E) (x0 : E)
        (a b : R) (N : Nat) (t : R) : E :=
      let h := (b - a) / N
      let k := min (Nat.floor ((t - a) / h)) (N - 1)
      let xk := eulerGrid v x0 a h k
      let tk := a + k * h
      xk + (t - tk) * v(tk, xk)
    \end{verbatim}
    The \texttt{min} with $N-1$ handles the right endpoint $t = b$.

  \item \textbf{Right derivatives at grid points.}
    At $t = t_k$ (a grid point), the interpolant has different
    left and right derivatives. The Gronwall theorem only requires
    the right derivative (\texttt{HasDerivWithinAt} with
    \texttt{Set.Ici t}), which matches the derivative of the
    piece on $\Icc{t_k}{t_{k+1}}$.

  \item \textbf{Extending $v$ beyond $\Icc{a}{b}$.}
    The Gronwall theorem \texttt{dist\_le\_of\_approx\_trajectories\_ODE}
    requires $v(t, \cdot)$ to be $K$-Lipschitz for \emph{all}
    $t \in \R$, not just $t \in \Icc{a}{b}$. If $v$ is only
    defined on $\Icc{a}{b}$, extend it by setting
    $v(t, x) = v(a, x)$ for $t < a$ and $v(t, x) = v(b, x)$ for
    $t > b$. Alternatively, use the variant
    \texttt{dist\_le\_of\_approx\_trajectories\_ODE\_of\_mem}
    with the time-dependent set $s(t) = E$.

  \item \textbf{NNReal for constants.}
    Following Mathlib conventions, consider using $\R_{\ge 0}$
    (\texttt{NNReal}) for $K$, $L$, $M$, $h$ to avoid carrying
    non-negativity proofs.

  \item \textbf{The Picard--Lindel\"of hypotheses.}
    The structure \texttt{IsPicardLindelof} requires constants
    as \texttt{NNReal} and the time parameter $t_0$ as a subtype
    of $\Icc{t_{\min}}{t_{\max}}$. Set $t_0 = a$,
    $t_{\min} = a$, $t_{\max} = b$, $r = 0$,
    $a_{\mathrm{ball}} = M \cdot (b - a)$ (as an NNReal, requiring
    $M \ge 0$ and $b \ge a$), $L_{\mathrm{bound}} = M$, and
    $K_{\mathrm{lip}} = K$.
\end{enumerate}


\end{document}
